\chapter{Discussion}

\section{Overview of Key Findings}

This chapter provides a critical interpretation of the experimental results presented in Chapter 4, examining the performance characteristics of traditional machine learning algorithms for partial discharge classification and the influence of feature selection strategies on diagnostic accuracy \parencite{Aldosari_2023, Govindarajan_2023_1}. The discussion addresses three fundamental aspects of the research: the comparative evaluation of machine learning algorithms under standardized experimental conditions, the identification and interpretation of optimal feature subsets through independent selection regimes, and the alignment of observed performance patterns with established literature in partial discharge diagnostics \parencite{Srivastava_2023_1, Zemouri_2023_1}. The analysis focuses on understanding the underlying mechanisms driving classification performance rather than merely reporting numerical outcomes, with particular attention to the trade-offs between different performance metrics and the physical interpretability of selected feature representations \parencite{Seitz_2024_1, Zhang_2024_29}. This interpretive framework enables evidence-based conclusions regarding algorithm suitability for field deployment scenarios and provides insights into the stability and diagnostic relevance of multi-domain feature engineering approaches in partial discharge monitoring applications \parencite{Long_2024, Florkowski_2024}.

\section{Discussion Based on Research Objectives}

\subsection{Comparative Evaluation of Traditional Machine Learning Algorithms}

The experimental results establish a clear performance hierarchy among the evaluated machine learning classifiers when applied to partial discharge classification under both feature selection regimes \parencite{Govindarajan_2023_1, Aldosari_2023}. Under Track 4C (MLJAR feature selection), Random Forest achieved the highest F1 score of 73.11\% alongside the highest ROC-AUC of 90.63\%, with an accuracy of 91.83\% and precision of 81.50\%. Under Track 4B (Featurewiz selection), Random Forest similarly achieved the highest F1 score of 71.84\% with a ROC-AUC of 90.24\%, accuracy of 91.91\%, and precision of 83.55\%. Artificial Neural Networks demonstrated the highest recall performance across both tracks, achieving 72.32\% on Track 4C and 69.21\% on Track 4B, with corresponding F1 scores of 72.65\% and 70.23\% respectively \parencite{Srivastava_2023_1, Aldosari_2023}. Support Vector Machines achieved the highest precision values, reaching 84.50\% on Track 4B and 84.49\% on Track 4C, with accuracy values of 91.64\% and 91.55\% respectively, and F1 scores of 69.13\% and 68.54\% \parencite{Govindarajan_2021_3, Aldosari_2023}.

The superior performance of Random Forest classifiers across both feature selection tracks indicates that ensemble methods provide optimal balance between precision and recall for partial discharge classification tasks \parencite{Govindarajan_2023_1, Aldosari_2023}. This outcome demonstrates that the complex, non-linear relationships inherent in partial discharge signals are more effectively captured through ensemble averaging of multiple decision trees rather than through single-tree models or linear margin-based classifiers \parencite{Govindarajan_2021_3, Srivastava_2023_1}. The ability of Random Forest to maintain robust performance across substantially different feature set sizes (31 features versus 2070 features) suggests that its ensemble mechanism provides inherent regularization that mitigates overfitting risks associated with high-dimensional feature spaces \parencite{Aldosari_2023, Zemouri_2023_1}. This finding validates the research objective of identifying the most suitable machine learning algorithm for partial discharge classification, as Random Forest consistently outperformed alternative approaches while maintaining interpretability through feature importance rankings \parencite{Govindarajan_2023_1, Seitz_2024_1}. The observed performance advantage of Artificial Neural Networks in recall metrics, despite lower overall F1 scores compared to Random Forest, indicates that neural networks can effectively exploit higher-order feature interactions and complex decision boundaries that enable improved detection sensitivity for minority class instances \parencite{Srivastava_2023_1, Aldosari_2023}. However, the superior balanced performance of Random Forest, as evidenced by higher F1 scores, suggests that ensemble methods provide a more optimal trade-off between detection sensitivity and prediction confidence \parencite{Govindarajan_2023_1, Zemouri_2023_1}.

The observed dominance of Random Forest in the present study aligns with its widespread adoption in partial discharge pattern recognition research utilizing engineered feature representations \parencite{Govindarajan_2023_1, Aldosari_2023}. Previous studies have similarly reported strong performance of Random Forest classifiers in PD diagnostics, with applications ranging from switchgear fault diagnosis to cable insulation assessment \parencite{Govindarajan_2021_3, Seitz_2024_1}. The consistency between the present results and prior literature suggests that ensemble tree methods represent a robust algorithmic choice for partial discharge classification tasks, particularly when working with multi-domain feature representations extracted from field monitoring data \parencite{Aldosari_2023, Long_2024}. The strong performance of Support Vector Machines observed in this study, achieving precision values exceeding 84\% across both tracks, is consistent with their frequent application in PD source classification and defect type recognition \parencite{Govindarajan_2021_3, Aldosari_2023}. Previous research has demonstrated SVM effectiveness in high-dimensional feature spaces, which aligns with the present findings where SVM maintained competitive performance despite the substantial dimensionality differences between Track 4B (31 features) and Track 4C (2070 features) \parencite{Govindarajan_2023_1, Zemouri_2023_1}. However, the present results indicate that while SVM provides excellent precision, Random Forest achieves superior balanced performance through higher F1 scores, suggesting that ensemble methods may be preferable when both detection sensitivity and prediction confidence are operationally important \parencite{Govindarajan_2023_1, Adhikari_2024}. The observed performance of Artificial Neural Networks, achieving the highest recall values but lower overall F1 scores compared to Random Forest, reflects a pattern that has been reported in previous PD diagnostics research \parencite{Srivastava_2023_1, Aldosari_2023}. Prior studies have noted that neural networks can effectively capture complex feature interactions but may require careful regularization and hyperparameter tuning to achieve optimal balanced performance \parencite{Aldosari_2023, Zemouri_2023_1}. The present results extend this understanding by demonstrating that while neural networks excel in recall-oriented applications, ensemble methods provide superior overall performance when multiple metrics must be simultaneously optimized \parencite{Govindarajan_2023_1, Srivastava_2023_1}.

\subsection{Identification of Optimal Feature Sets via Feature Selection}

The feature importance rankings obtained from the Featurewiz--XGBoost and MLJAR internal selection frameworks reveal an apparent divergence in the prioritization of individual descriptors. At first glance, classical statistical features such as kurtosis and skewness appear to be ranked lower than spectral entropy and composite interaction features. However, closer examination indicates strong conceptual agreement regarding the dominant signal characteristics governing partial discharge (PD) classification. This finding is important because it demonstrates that different feature selection mechanisms may emphasize distinct mathematical representations of the same underlying physical phenomena. Consequently, the observed differences should be interpreted not as inconsistency, but as evidence that PD-related information is encoded redundantly across multiple feature domains \parencite{Yao_2020_1, Du_2023}.

Previous studies on partial discharge (PD) signal classification have predominantly relied on handcrafted feature sets derived from time-domain statistics, frequency-domain summaries, and phase-resolved partial discharge (PRPD) representations. In particular, time-domain statistical descriptors such as kurtosis, skewness, and crest factor are frequently reported as dominant features due to their sensitivity to impulsive discharge events and asymmetric amplitude distributions. Frequency-domain features, including FFT-based energy measures and spectral peaks, are also commonly employed to characterize discharge-related spectral content, while PRPD-based quadrant statistics are widely used to encode phase-locked discharge behaviour. In most prior works, these features are evaluated either individually or within relatively small, manually curated feature sets, with limited use of interaction or composite descriptors \parencite{Auni_2020, Soh_2022}.

Under the Featurewiz--XGBoost framework, spectral entropy and related frequency-domain interaction features dominate the ranking, while classical statistical descriptors appear with lower relative importance. This outcome reflects the sensitivity of ANOVA-based ranking to features that maximize class separability in the spectral domain. The prominence of spectral entropy indicates that PD signals exhibit structured, non-random frequency distributions that are distinct from background noise, confirming that discharge activity introduces measurable disorder and localized energy concentration in the frequency spectrum \parencite{Sun_2020, Romphuchaiyapruek_2025}. This result shows that PD events are not solely characterized by impulsiveness in the time domain, but also by organized patterns in spectral space.

In contrast, the MLJAR internal feature selection framework prioritizes composite and interaction features that combine spectral entropy with fractal dimension, RMS, variance, and PRPD-related measures. The dominance of these hybrid descriptors indicates that PD classification benefits from capturing joint relationships between signal complexity, amplitude modulation, and phase-resolved behaviour. This suggests that PD signals exhibit multi-scale structure, where impulsiveness, complexity, and energy distribution interact in a non-linear manner. Importantly, many of the highest-ranked MLJAR features implicitly encode impulsiveness and asymmetry through composite formulations, even when kurtosis or skewness do not appear explicitly among the top-ranked variables.

When compared with previous studies, the apparent reduction in ranking of classical time-domain statistics can be explained primarily by differences in feature construction and selection methodology rather than by disagreement in physical interpretation. Earlier works typically evaluate kurtosis, skewness, and crest factor as standalone descriptors within limited feature spaces, causing these features to emerge as dominant indicators of PD activity. In contrast, the present study employs automated feature engineering and independent, interaction-aware selection frameworks, which redistribute importance toward composite features that subsume impulsiveness and asymmetry within higher-order representations. As a result, classical statistical features may appear ranked lower, not because they are less informative, but because their diagnostic contribution is absorbed into interaction features that capture the same physical properties more comprehensively.

Overall, this finding refines rather than contradicts existing literature. It demonstrates that the perceived importance of individual PD features is strongly influenced by the feature selection paradigm. When interaction-aware and model-agnostic selection is applied, PD discrimination is driven not by isolated descriptors, but by coupled representations of complexity, spectral organization, and impulsiveness. Thus, the present results extend previous studies by showing that impulsive and asymmetric discharge behaviour remains central to PD characterization, whether expressed explicitly through traditional time-domain statistics or implicitly through composite spectral and structural features.



% \section{Interpretation of Model Performance and Feature Characteristics}

% A fundamental pattern emerging from the experimental results is the clear separation between overall classification correctness (accuracy) and balanced error trade-offs (F1 score), which reflects the inherent challenges of imbalanced diagnostic datasets \parencite{Zemouri_2023_1, Srivastava_2023_1}. This behaviour is consistent with well-established limitations of accuracy as a performance metric under class imbalance, where models can achieve high accuracy by predominantly predicting the majority class while failing to detect critical minority class instances \parencite{Zemouri_2023_1, Adhikari_2024}. In the present study, despite achieving accuracy values exceeding 87\% across all models, recall values for the best-performing models ranged from 67.09\% to 72.32\%, indicating that while performance is substantially improved compared to baseline methods, a proportion of actual fault instances remain undetected.

% For PD monitoring applications, recall-oriented performance interpretation is often emphasized because missed detections (false negatives) can have severe risk implications in condition monitoring contexts, potentially leading to progressive insulation degradation and unexpected equipment failures \parencite{Zemouri_2023_1, Adhikari_2024}. The observed precision--recall trade-offs, where Random Forest achieved the highest precision (83.55\% on Track 4B, 81.50\% on Track 4C) with strong recall (67.09\% on Track 4B, 68.94\% on Track 4C), versus ANN achieving the best recall (69.21\% on Track 4B, 72.32\% on Track 4C) with moderate precision (71.61\% on Track 4B, 73.41\% on Track 4C), illustrate the fundamental tension between minimizing false alarms and maximizing fault detection sensitivity \parencite{Zemouri_2023_1, Adhikari_2024}.

% The persistence of similar jointly selected features across both feature selection tracks provides strong evidence for the stability and physical interpretability of the engineered descriptors \parencite{Seitz_2024_1, Zhang_2024_29}. Time--frequency descriptors and entropy-like measures (spectral entropy, Shannon entropy combinations) were consistently prioritized, which aligns with their established role in reflecting spectral complexity and signal irregularity under noise and interference conditions \parencite{Bao_2024_1, Friebe_2024}. Phase-resolved PRPD summaries, particularly quadrant-based differences and polarity ratios, were extensively retained, supporting their interpretation as proxies for discharge phase preference and polarity asymmetry under phase-locked AC excitation \parencite{Orellana_2023_1, Seitz_2024_1}.

% The strong balanced performance of neural network classifiers (ANN achieving F1 scores of 72.65\% and 70.23\% on Tracks 4C and 4B, respectively) can be attributed to their ability to exploit higher-order feature interactions and nonlinear decision boundaries in the engineered representations \parencite{Aldosari_2023, Srivastava_2023_1}. However, Random Forest achieved superior F1 scores on both tracks (73.11\% on Track 4C, 71.84\% on Track 4B), indicating that ensemble methods provide the best overall balanced performance \parencite{Govindarajan_2023_1, Aldosari_2023}. Such performance advantages are typically conditional on careful hyperparameter tuning, regularization strategies, and validation protocols in noisy monitoring settings, as evidenced by the standard deviations observed in precision and recall metrics \parencite{Srivastava_2023_1, Zemouri_2023_1}.

% The strong discrimination capability demonstrated by Random Forest classifiers, achieving the highest ROC-AUC (90.63\%) and F1 score (73.11\%) on Track 4C, is consistent with their frequent application in PD diagnostics with high-dimensional feature spaces, where their ensemble averaging provides robust generalization \parencite{Govindarajan_2021_3, Aldosari_2023}. SVM classifiers also demonstrated strong performance, achieving high precision (84.50\% on Track 4B, 84.49\% on Track 4C) and accuracy (91.64\% on Track 4B, 91.55\% on Track 4C), which can be attributed to their margin-maximization objective in high-dimensional feature spaces \parencite{Govindarajan_2021_3, Aldosari_2023}. Ensemble tree models (Random Forest) exhibited balanced precision--recall behaviour, achieving the highest precision values (83.55\% on Track 4B, 81.50\% on Track 4C) alongside strong recall (67.09\% on Track 4B, 68.94\% on Track 4C), which may be interpreted as providing optimal trade-offs between high-confidence predictions and detection sensitivity \parencite{Zemouri_2023_1, Adhikari_2024}.

% \section{Comparison with Existing Literature}

% The benchmark comparison presented in Chapter 4 must be interpreted with appropriate caution, as reported performance metrics across different studies are often not directly comparable due to fundamental differences in experimental design, data characteristics, and evaluation protocols \parencite{Aldosari_2023, Sikorski_2025}. Differences in sensing modality (contactless antenna versus galvanic coupling, HFCT, or UHF sensors), defect taxonomy and classification granularity, and evaluation methodology (cross-validation strategy, class balance, label quality) are widely cited as primary drivers of apparent performance discrepancies in the PD diagnostics literature \parencite{Vera_2023_1, Long_2024}.

% Studies emphasizing cable terminal PD signals, GIS-specific discharge patterns, or PRPD-image-based classification pipelines may leverage data structures and feature representations that differ substantially from the antenna-based long-term monitoring signals utilized in the present work \parencite{Aldosari_2023, Long_2024}. Moreover, studies reporting high accuracy values (ranging from 91.5\% to 97.0\% in the cited literature) may differ materially in class balance, label quality and annotation consistency, noise cancellation assumptions, and the specific fault types included in the classification task \parencite{Zemouri_2023_1, Friebe_2024}. The present study's best accuracy of 91.91\% (Random Forest on Track 4B) and 91.83\% (Random Forest on Track 4C) positions these results within the lower-to-mid range of reported values, which is reasonable given the realistic field monitoring context and class imbalance challenges \parencite{Florkowski_2024, Sikorski_2025}.

% Accordingly, the present results are most appropriately positioned as an interpretable baseline within a multi-domain feature engineering framework and conventional machine learning model zoo, rather than as a direct performance comparison against studies utilizing fundamentally different data acquisition and processing methodologies \parencite{Aldosari_2023, Srivastava_2023_1}. The continued recommendation to report multiple complementary performance metrics (accuracy, precision, recall, F1 score, ROC-AUC) is consistent with established diagnostic evaluation guidance in imbalanced classification settings, where single-metric optimization can lead to misleading conclusions about practical utility \parencite{Zemouri_2023_1, Adhikari_2024}.

\section{Strengths of the Proposed Approach}

A principal strength of this study lies in the utilization of long-term, multi-site field monitoring data acquired from operational European overhead power lines, which directly supports realism and deployment relevance in partial discharge (PD) analytics \parencite{Florkowski_2024, Sikorski_2025}. 
Unlike laboratory-based datasets, the use of field-acquired measurements enables the evaluation of model performance under authentic operating conditions, capturing signal variability arising from environmental interference, load dynamics, and insulation aging processes \parencite{Orellana_2023_1, Florkowski_2024}. 
The multi-year temporal coverage and geographic diversity of the dataset further allow for the assessment of model stability and consistency across heterogeneous real-world conditions \parencite{Seitz_2024_1, Sikorski_2025}.

Another key strength is the adoption of a comprehensive multi-domain feature engineering framework encompassing time-domain statistics, frequency-domain spectral descriptors, time--frequency representations, and phase-resolved partial discharge (PRPD) features \parencite{Aldosari_2023, Long_2024}. 
This design explicitly acknowledges that different PD mechanisms manifest complementary characteristics across multiple signal domains, thereby improving class separability compared to single-domain feature representations \parencite{Aldosari_2023, Long_2024}. 
By integrating physically meaningful descriptors linked to discharge energy, temporal behavior, spectral content, and phase dependency, the proposed feature set provides a robust and interpretable representation of PD phenomena \parencite{Florkowski_2024, Long_2024}.

The use of two independent feature selection pipelines, namely Featurewiz and MLJAR, constitutes an additional methodological strength by enabling cross-validation of feature relevance under distinct selection philosophies \parencite{Seitz_2024_1, Zhang_2024_29}. 
The substantial overlap observed between the features selected by both methods provides internal evidence of selection stability and reduces the risk of algorithm-specific bias in high-dimensional feature spaces \parencite{Zhang_2024_29, Seitz_2024_1}. 
This convergence strengthens confidence that the retained features capture physically meaningful and diagnostically relevant characteristics rather than dataset-specific artifacts \parencite{Seitz_2024_1, Long_2024}.

Furthermore, the evaluation framework employs multiple complementary performance metrics, including accuracy, precision, recall, F1-score, and ROC-AUC, enabling a risk-aware interpretation of model behavior under class imbalance \parencite{Zemouri_2023_1, Adhikari_2024}. 
Such comprehensive reporting supports application-driven assessment, where different operational contexts may prioritize fault sensitivity, false-alarm reduction, or overall classification robustness \parencite{Zemouri_2023_1, Adhikari_2024}. 
By pairing quantitative performance metrics with uncertainty estimates derived from stratified cross-validation, the proposed approach facilitates informed decision-making for industrial PD monitoring deployments \parencite{Seitz_2024_1, Florkowski_2024}.


